\section{Sacrifice, Visible and Invisible}

Sacrifice is the act of religion with the strangest résumé. Every other religious act has skeptical explanations ready at hand; sacrifice has mostly embarrassment. Killing an animal, burning grain, pouring out wine---destroying good things on purpose, for no one's use---looks, to a certain modern eye, like the clearest possible case of religion as waste. Yet the practice is as close to universal as anything in the human record, and Aquinas opens his question on it from exactly that datum: ``At all times and among all nations there has always been the offering of sacrifices.''\endnote{\emph{Summa theologiae} II-II, q.~85, a.~1, \latin{sed contra}; checked 2026-07-24, English at New Advent, Latin at Corpus Thomisticum. The argument-form---what all peoples do is seemingly natural---is an appeal to natural inclination, not a statistical survey; the corpus supplies the reason the inclination exists.} Whatever sacrifice is, it is not a local invention. Something in the creature as such produces it, the way something in the creature produces language.

Aquinas's account of that something is a compressed anthropology. Natural reason tells a man that he is subject to something above him: he registers ``the defects which he perceives in himself, and in which he needs help and direction from someone above him''---he did not make himself, does not sustain himself, cannot complete himself. Natural reason likewise dictates that submission and honor be shown to what is above one. But---here is the human twist---``the mode befitting to man is that he should employ sensible signs in order to signify anything, because he derives his knowledge from sensibles.'' A pure spirit could submit in pure act; man must \emph{say} it with things. ``Hence it is a dictate of natural reason that man should use certain sensibles, by offering them to God in sign of the subjection and honor due to Him, like those who make certain offerings to their lord in recognition of his authority.''\endnote{\emph{Summa theologiae} II-II, q.~85, a.~1, corpus (\latin{ex naturali ratione procedit quod homo quibusdam sensibilibus rebus utatur, offerens eas Deo, in signum debitae subiectionis et honoris}); checked 2026-07-24 as above. The reply to the first objection adds the boundary this essay observes: the offering of sacrifice belongs \emph{generically} to the natural law, ``but the determination of sacrifices is established by God or by man''---which victims, rites, and times is positive determination, not natural deduction. No claim is made here that any particular rite is naturally required.} Sacrifice, in its natural-law core, is the tenant's basket of first fruits carried to a landlord who owns the orchard, the basket, and the tenant: materially absurd, and precisely therefore perfectly clear as a \emph{sign}. Nothing is transferred; everything is acknowledged. Which things, which rites, which days---all that, Aquinas immediately adds, is determination, positive institution, not natural deduction. Nature dictates that the acknowledgment be enacted in stuff; it does not write the rubrics.

What the enacted sign signifies, he states with a precision that dissolves the embarrassment: ``the sacrifice that is offered outwardly represents the inward spiritual sacrifice, whereby the soul offers itself to God,'' and it offers itself to God ``as its beginning by creation, and its end by beatification.''\endnote{\emph{Summa theologiae} II-II, q.~85, a.~2, corpus, citing Psalm 50:19 (``A sacrifice to God is an afflicted spirit'') and Augustine, \emph{De civitate Dei} X.19 on words offered to God signifying what the heart offers; checked 2026-07-24 as above. The corpus draws the exclusivity conclusion: because the soul offers itself to God alone as beginning and beatifying end, outward sacrifice too is offered to God alone---the act of latria that cannot be transferred to any creature.} Beginning by creation; end by beatification: the two anchors of Part I's whole argument, reappearing as the content of the burnt offering. The animal on the altar was never the payment; it was the pronoun. \emph{This---this life, this substance, this mine---stands for me; I am the thing being handed over.} And because the self being handed over is received-being all the way down, the handing-over adds nothing to God and was never supposed to: it is the acknowledgment, in the costliest available sign, that there was never anything to withhold.

The deepest ancient meditation on this grammar is Augustine's, in the tenth book of \emph{The City of God}, and it repays quotation because it contains, at fifteen centuries' distance, both of this series' governing images. God needs nothing, Augustine begins: ``who is so foolish as to suppose that the things offered to God are needed by Him for some uses of His own?'' Right worship ``profits not Him, but man. For no man would say he did a benefit to a fountain by drinking, or to the light by seeing.''\endnote{Augustine, \emph{De civitate Dei} X.5, in the public-domain translation of Marcus Dods (\emph{The City of God}, vol.~I, in \emph{The Works of Aurelius Augustine, A New Translation}, Edinburgh: T.\,\&\,T. Clark, 1871; Book X is Dods's own translation within the collaborative edition); wording checked 2026-07-24 in the source library's tracked transcription of the 1871 printing, at the registered passage covering X.5--6, and against the New Advent transcription of the same translation. The fountain and the light are Augustine's own vehicles for the point Part I made with the spring and the lamp; the convergence is inherited from him, not engineered by this essay.} The fountain and the light: Part I's spring and lamp, in Augustine's own hand, deployed for exactly this point---the drinker and the seer receive everything and confer nothing. Then the definition that governs everything after it: ``A sacrifice, therefore, is the visible sacrament or sacred sign of an invisible sacrifice.'' And then, in the next chapter, the widening: ``Thus a true sacrifice is every work which is done that we may be united to God in holy fellowship, and which has a reference to that supreme good and end in which alone we can be truly blessed.'' Mercy shown to a neighbor, the body disciplined, the soul turned---each, done for God's sake, is sacrifice in the proper sense; the visible rites were always signs of this.\endnote{\emph{De civitate Dei} X.5 (the visible sacrament or sacred sign) and X.6 (the true-sacrifice definition; mercy for God's sake as sacrifice; the body and the soul as offerings, citing Romans 12:1--2), Dods translation, checked 2026-07-24 as above. Augustine's X.6 argument runs through Romans 12:1---``I beseech you therefore, brethren, by the mercy of God, that you present your bodies a living sacrifice, holy, pleasing unto God, your reasonable service'' (Douay--Rheims, checked 2026-07-24)---and its ``living sacrifice'' is the hinge on which the chapter turns from definition to ecclesiology. Note the continuity with II-II, q.~81, a.~1, ad~1 (section~3 above): what Aquinas parses as religion \emph{commanding} the act of mercy, Augustine states as mercy \emph{being} true sacrifice when referred to God.}

To this point, the argument has stayed on the ascent's side of the seam this series keeps visible: everything said of sacrifice so far---its natural-law root, its sign-grammar, its terminus in the self-offering of the soul---is available to reason reflecting on creaturehood. But Augustine does not stop there, and honesty requires marking where his next sentence stands. The chapter's climax is not a definition but a confession: the ``whole redeemed city, that is to say, the congregation or community of the saints, is offered to God as our sacrifice through the great High Priest, who offered Himself to God in His passion for us, that we might be members of this glorious head''; and so, he concludes, ``This is the sacrifice of Christians: we, being many, are one body in Christ''---the sacrifice the Church ``continually celebrates in the sacrament of the altar,\ldots{} in which she teaches that she herself is offered in the offering she makes to God.''\endnote{\emph{De civitate Dei} X.6, Dods translation, checked 2026-07-24 as above; the interior quotation is Romans 12:5. These sentences are cited as Christian confession and the Church's teaching, not as conclusions of the natural-theological thread: Christ's unique saving oblation is prior, and creatures are incorporated into it rather than left to duplicate it. The philosophical argument of this essay establishes that acknowledgment in sign is owed; it does not and cannot establish the Priest, the Passion, or the altar.} Philosophy can demonstrate that an offering is owed. It cannot discover the Cross or manufacture the Eucharist from creaturely dependence. What revelation claims---and it is a claim of faith, not a theorem---is that the incarnate Son freely offered his human life to the Father for our salvation, and that creatures are invited to be \emph{included in} his prior oblation as members of the body being offered.

And with that, the psalmist's answer finally parses. Asked what he will render for all that has been rendered to him, he does not reach for an asset; he reaches for the cup: ``I will take the chalice of salvation; and I will call upon the name of the Lord.'' \emph{Taking} is the rendering. In a relation with no return current, where every payment would be made in the creditor's own coin, the one thing the debtor can genuinely contribute is the manner of his receiving: to take the gift knowingly, name the giver aloud, and bind himself publicly to go on doing so---``I will pay my vows to the Lord before all his people.'' The Latin verbs face each other across the verse like a diptych: \latin{quid retribuam}---\latin{accipiam}. What shall I give back? I shall receive. It is not evasion; it is precision. Reception with the eyes open is the only currency that was ever really the creature's own, and the psalm, the analysis of gratitude, and the theology of sacrifice all converge on it from their three directions.
